\documentclass[10pt, letterpaper]{paper}
\usepackage[round]{natbib}
\setcitestyle{authoryear}
\usepackage[margin=2.5cm]{geometry}
\usepackage{amsmath,amssymb,amsthm}
\usepackage{booktabs}
\usepackage{graphicx}
\usepackage[colorlinks=true,linkcolor=blue,citecolor=blue]{hyperref}
\usepackage{doi}
\usepackage{orcidlink}
\usepackage{authblk}
\renewcommand{\Affilfont}{\footnotesize}
\renewcommand{\Authfont}{\normalsize}

\newtheorem{theorem}{Theorem}[section]
\newtheorem{proposition}[theorem]{Proposition}
\newtheorem{lemma}[theorem]{Lemma}
\newtheorem{corollary}[theorem]{Corollary}
\newtheorem{definition}[theorem]{Definition}
\theoremstyle{remark}
\newtheorem{remark}[theorem]{Remark}
\newtheorem{principle}[theorem]{Principle}

\newcommand{\R}{\mathbb{R}}
\newcommand{\E}{\mathbb{E}}
\newcommand{\eqd}{\overset{d}{=}}
\newcommand{\Phifun}{\Phi}
\newcommand{\Amed}{\mathcal{A}_{\mathrm{med}}}

\title{\textbf{Functional detrended fluctuation analysis (FDFA):\\
working notes --- the idea, the axioms, and the
prediction--falsification record}}

\author{D. Sierra-Porta\orcidlink{0000-0003-3461-1347}}
\affil{Universidad Tecnológica de Bolívar. Grupo de Investigación Gravitación y Matemática Aplicada - GIGMA \& Grupo de Investigación Física Aplicada y Procesamiento de Imágenes y Señales - FAPIS., Parque Industrial y Tecnológico Carlos Vélez Pombo Km 1 Vía Turbaco, Cartagena de Indias, 130010, Bolívar, Colombia}
\affil{Corresponding author: dporta@utb.edu.co}
\date{\today}

\begin{document}
\maketitle

\begin{abstract}
\noindent These are the working notes behind the FDFA project. They
record, in the order in which it happened, the question that started
the project, the minimal theory built to make the question precise,
the falsifiable predictions derived from that theory, and the outcome
of each prediction --- including one falsification that became the
most useful structural insight of the program. The notes accompany
the code repository that implements everything described here; the
consolidated manuscript is a separate document. Nothing in these
notes is required to use the code, but they explain \emph{why} the
code is organized the way it is.
\end{abstract}

\section{The question}

Classical DFA/MFDFA \citep{peng1994,kantelhardt2002} computes, in
each window $v$ of scale $s$, the \emph{variance} of the detrending
residuals,
\begin{equation}
F^2(v,s)=\frac1s\sum_{i=1}^{s}
\bigl(Y((v-1)s+i)-P_m^{(v)}(i)\bigr)^2 ,
\label{eq:classicF}
\end{equation}
and aggregates the windows with $q$-th power means,
$F_q(s)=(\tfrac1{2N_s}\sum_v[F^2(v,s)]^{q/2})^{1/q}\sim s^{h(q)}$.
Essentially every ingredient of this pipeline has been varied in the
literature --- the detrending, the windowing, the correlation
structure probed --- except the fluctuation measure itself, which has
remained the variance in every variant. The project starts from one
question: \emph{is the variance necessary, and what exactly would
change if it were replaced?}

The answer developed below: for the scaling exponent the variance is
irrelevant (any admissible dispersion functional gives the same
exponent); what the choice controls is the statistical behavior of
the estimator; and the decisive design choice turns out to sit one
level up --- in how windows are aggregated --- which in turn yields a
surrogate-free diagnostic of the origin (correlational versus
distributional) of multifractality.

\section{The construction}

\begin{definition}[Dispersion functional]\label{def:phi}
$\Phifun:\bigcup_{s\ge1}\R^s\to[0,\infty)$, measurable on each
$\R^s$, satisfying
(A1) \emph{positive homogeneity}, $\Phifun(\lambda r)=\lambda\Phifun(r)$
for $\lambda>0$;
(A2) \emph{symmetry}, $\Phifun(-r)=\Phifun(r)$, and permutation
invariance;
(A3) \emph{nondegeneracy and continuity};
(A4) \emph{sampling consistency}: for every c\`adl\`ag
$f:[0,1]\to\R$, $\Phifun(f(1/s),\dots,f(s/s))$ converges as
$s\to\infty$ to a nondegenerate limit $\Phifun_\infty(f)$
(c\`adl\`ag, not merely continuous, so that jump processes are
covered).
\end{definition}

\begin{definition}[FDFA, two layers]\label{def:fdfa}
$F_\Phifun(v,s)=\Phifun(r^{(v,s)})$ on the window residuals (inner
layer), then
$F_{\Phifun,\mathcal A}(s)=\mathcal A(\{F_\Phifun(v,s)\}_v)$ across
windows (outer layer), and the exponent from the log-log slope.
Members are written $\mathrm{FDFA}_{\Phifun,\mathcal A}$. Classical
DFA is $\mathrm{FDFA}_{\mathrm{RMS},\mathcal A_2}$; MFDFA is the
slice $\{\mathrm{FDFA}_{\mathrm{RMS},\mathcal A_q}\}_q$.
\end{definition}

The organizing idea: \textbf{the inner functional decides what
``fluctuation'' means inside a window; the outer aggregator decides
which windows dominate.} Implemented inner functionals (all
normalized to Gaussian $\sigma$-consistency): RMS, $L^p$
($p=1,\,1/2$), MAD, quantiles $Q_{0.75}$, $Q_{0.90}$, IQR, Huber
scale, trimmed standard deviation, $L^\infty$, entropy power
$e^{\hat h}$. Analyzed and excluded: bounded functionals (kernel
MMD, total variation, KL --- no scaling law is possible),
scale-invariant functionals, and reducible ones
($W_p(\cdot,\delta_0)$ is exactly $L^p$; energy distance is
$L^1$-type). Entropies are admissible only after exponentiation.

\section{What the theory says}

\begin{lemma}[Dilation equivariance]\label{lem:commute}
For $\Pi_s$ the $L^2[0,s]$-projection onto polynomials of degree
$\le m$ and $(D_af)(t)=f(at)$:
$(I-\Pi_s)g=D_{1/s}(I-\Pi_1)D_sg$, and $(I-\Pi_s)$ commutes with
positive scalars.
\end{lemma}

\begin{proposition}[Exponent universality]\label{prop:universal}
Let $X$ be $H$-sssi with c\`adl\`ag paths, $\Phifun$ admissible,
and let $Z_\Phifun=\Phifun_\infty((I-\Pi_1)X|_{[0,1]})$ have finite
nonzero median. Then $s^{-H}F_\Phifun(v,s)\to_dZ_\Phifun$, hence
$\operatorname{med}F_\Phifun(v,s)=C_\Phifun s^H(1+o(1))$, and the
same for the mean under uniform integrability. The exponent is $H$
for \emph{every} admissible $\Phifun$.
\end{proposition}

The proof is three lines given the lemma: residuals of the
$s$-window are, in distribution, $s^H$ times time-rescaled residuals
of the unit window; homogeneity extracts $s^H$; sampling consistency
gives the limit. Consequence: replacing the variance never changes
the estimated quantity, only the estimator's statistics --- the
\emph{license} for everything else.

\begin{remark}[Finite-range correction of $L^\infty$]\label{rem:linf}
The supremum is admissible, but the sampled maximum approaches the
continuum supremum from below with a slowly varying deficit (in the
iid caricature, a $\sqrt{2\log s}$ factor; dependence within the
window slows but does not remove it), inflating finite-range slopes
--- more at small $H$. Confirmed numerically; leader-type
normalization \citep{jaffard2004} is the cure.
\end{remark}

\begin{proposition}[Heavy tails]\label{prop:stable}
For iid symmetric $\alpha$-stable input, $0<\alpha<2$ (profile
$=1/\alpha$-self-similar stable walk): (i)
$s^{-1/\alpha}F_\Phifun(v,s)\to_d\Xi$ with regularly varying tail of
index $\alpha$; (ii) for the $q$-mean outer layer with $q>\alpha$,
$\E\Xi^q=\infty$ and the bifractal law $h(q)=1/q$ of
\citet{nakao2000} applies --- in particular $h(2)=1/2$ for every
$\alpha<2$: \emph{classical DFA is asymptotically blind to the
self-similarity index of stable noise}; (iii) a median outer layer
recovers $1/\alpha$ with any admissible inner functional.
\end{proposition}

\begin{principle}[Trend absorption]\label{pri:trend}
An extreme increment enters the profile as a step; after detrending
it is a perturbation of order $M$ spread over the whole window --- a
shape distortion, not a point outlier. No translation-insensitive
inner functional, robust or not, can reject it. Robustness against
heavy tails is effective only at the outer, cross-window layer.
\end{principle}

This principle was \emph{not} part of the original plan; it is the
post-mortem of a falsified prediction (P4 below), and it dictates
both the design rule (robustify the outer layer) and the diagnostic
that follows.

\section{The discriminator}

A $q$-mean with large $|q|$ is dominated by extreme windows, and a
window can be extreme for two indistinguishable reasons: genuinely
different local scaling (the multifractality one wants) or a
heavy-tailed event (distributional luck). Quantiles of the window
ensemble separate the two.

\begin{definition}[Quantile-exponent profile]\label{def:theta}
$\theta(\tau)$ is the log-log slope of the empirical
$\tau$-quantile of $\{F_\Phifun(v,s)\}_v$ against $s$, over scales
with $2N_s\gtrsim50$.
\end{definition}

If windows are exchangeable with a common scale factor (heavy-tailed
monofractal), every quantile scales identically: $\theta$ flat, even
when $\Delta h=h(-4)-h(4)$ is large. If window heterogeneity is
scale-structured (genuine multifractality), $\theta$ varies with
$\tau$ --- decreasing for cascades. Verdict rule, with thresholds
fixed on synthetic data before touching real data:
$\Delta h>0.15$ and $|\Delta\theta|\le0.08$ $\Rightarrow$
distributional; $\Delta h>0.15$ and $|\Delta\theta|>0.08$
$\Rightarrow$ genuine. This answers, from a single run and without
surrogate ensembles, the question customarily answered by shuffling
\citep{kantelhardt2002,movahed2006}; it is the fixed-$N$,
per-record counterpart of the asymptotic dichotomy of
\citet{kwapien2023}.

\section{The prediction ledger}

Every experiment tested a prediction formulated in advance. The
record, in full (details and figures in the repository):

\begin{table}[h]
\centering
\small
\begin{tabular}{clll}
\toprule
 & prediction & benchmark & outcome\\
\midrule
P1 & all admissible $\Phifun$ give slope $H$ & fGn,
$H\in\{0.3,0.5,0.7\}$ & \textbf{confirmed} (spread $<0.011$)\\
P2 & classical blind, robust recovers $1/\alpha$ & stable,
$\alpha\in\{1.2,1.5,1.8\}$ & \textbf{confirmed, strong form}\\
P3 & genuine width preserved by robust $\Phifun$ & binomial cascade
& \textbf{confirmed} ($0.23$--$0.25$ vs $0.215$)\\
P4 & robust inner $\Phifun$ suppresses spurious width & stable
$h(q)$ & \textbf{falsified} $\to$ Principle~\ref{pri:trend}\\
P5 & $(\Delta h,\Delta\theta)$ separates the three cases & fGn /
cascade / stable & \textbf{confirmed}\\
P6 & discriminator correct with Gaussian tails & MRW & \textbf{confirmed}
($\Delta\theta=-0.12$)\\
P7 & universality beyond fBm construction & ARFIMA$(0,d,0)$ &
\textbf{confirmed}\\
P8 & $\theta$-verdict $=$ shuffle verdict & five signal classes &
\textbf{confirmed} ($5/5$ at $n=2^{14}$)\\
P9 & RMSE $\propto n^{-1/2}$; order-$m$ bias floor & fGn, stable
grid & \textbf{confirmed} ($50$ real./point)\\
\bottomrule
\end{tabular}
\caption{The prediction--falsification record of the project.
P4 is reported with the same prominence as the confirmations: its
failure mechanism (trend absorption) reorganized the framework and
produced the discriminator.}
\label{tab:ledger}
\end{table}

Two caveats found along the way and kept: the discriminator needs
$n\gtrsim10^4$ under heavy tails (a stable record was misclassified
at $n=2^{12}$), and zero-inflated series (e.g.\ sunspot numbers at
solar minima) blow up negative-$q$ moments unless the lower scale
cut exceeds the zero-run scale --- with the default cut one obtains
a spurious $\Delta h\approx10$.

\section{Reality check (OMNI2)}

Applied to four heliophysical series with
$\mathrm{FDFA}^{(2)}_{\mathrm{Huber},\mathrm{med}}$ for exponents
and the discriminator for origin: sunspot $R$ reproduces the
surrogate-based conclusion of \citet{movahed2006} (width $0.62\to
0.01$ under shuffling; $\Delta\theta=-0.16$) from a single run;
$D_{st}$ and IMF $B$ come out genuinely multifractal, with visible
classical/robust exponent gaps flagging tail influence; and the IMF
increments $\delta B$ are a \emph{mixed} case --- shuffling retains
$\Delta h_{\mathrm{shuf}}=0.15$ of $0.38$, so roughly $40\%$ of the
width is distributional, refining the picture of
\citet{macek2011}.

\section{Open problems}

(1) Median outer layer sacrifices the Legendre formalism: robust
exponents or spectra, not both from one pipeline. (2)
$\Delta\theta$ needs a calibrated significance scale and a formal
three-way genuine/mixed/distributional verdict. (3) Rigor gaps:
discretization in Proposition~\ref{prop:universal}; tail constants
in Proposition~\ref{prop:stable}(i); median concentration under
long-range dependent window ensembles; quantitative trend
absorption; multiplicative-invariance version for cascades. (4)
Dictionary to the $p$-leader formalism \citep{leonarduzzi2016}.

\section*{Code}

Everything above is implemented in the accompanying repository:
\texttt{fdfa.py} (the two-layer estimator with pluggable
functionals), \texttt{generators.py} (benchmarks with known ground
truth), \texttt{run\_validation.py} (the nine-experiment ledger of
Table~\ref{tab:ledger}), \texttt{run\_application.py} (the OMNI2
analysis). All figures and tables regenerate from two seeded
commands in under one minute.

\bibliographystyle{plainnat}
\bibliography{refs}

\end{document}
